% M2 proof module: WP02 / WP03 / WP06 finite-radius resource laws.

\subsection{Finite-radius arbitrary-POVM survival laws}
\label{supp:finite-radius-full}

We use the convention
\begin{equation}
D_c=(A+A^\dagger)/2,
\qquad
D_s=(A-A^\dagger)/(2i).
\end{equation}
Let $P=\supp\rho_0$ and suppose the affine disk
\begin{equation}
\rho_0+xD_c+yD_s\succeq0
\qquad(x^2+y^2\le R^2)
\end{equation}
is physical for some $R>0$. Then positivity in both signs implies $A=PAP$.

Define on $P$
\begin{equation}
B=\rho_0^{-1/2}A\rho_0^{-1/2}.
\end{equation}
Writing $\beta=x-iy$, the affine family is congruent to
\begin{equation}
\id+\frac12(\beta B+\beta^*B^\dagger).
\end{equation}
For a unit vector $|\psi\rangle$, minimizing the expectation value over the phase of $\beta$ gives
\begin{equation}
1-|\beta|\,|\langle\psi|B|\psi\rangle|.
\end{equation}
Therefore the exact affine radius is
\begin{equation}
R_{\rm lin}=\frac1{w(B)}.
\label{eq:proof-Rlin}
\end{equation}
The standard numerical-radius inequality $\|B\|\le2w(B)$ gives
\begin{equation}
B B^\dagger\preceq\frac4{R_{\rm lin}^2}\id_P.
\end{equation}
By congruence,
\begin{equation}
A\rho_0^+A^\dagger
=\rho_0^{1/2}BB^\dagger\rho_0^{1/2}
\preceq\frac4{R_{\rm lin}^2}\rho_0.
\label{eq:proof-operator-majorization}
\end{equation}
Similarly,
\begin{equation}
A^\dagger\rho_0^+A
\preceq\frac4{R_{\rm lin}^2}\rho_0.
\end{equation}

For any POVM $\{M_y\}$ define
\begin{equation}
p_y=\Tr(\rho_0M_y),
\qquad
z_y=\Tr(AM_y).
\end{equation}
Weighted Hilbert--Schmidt Cauchy--Schwarz gives
\begin{equation}
\frac{|z_y|^2}{p_y}
\le\Tr(M_yA\rho_0^+A^\dagger).
\end{equation}
After summing the outcomes,
\begin{equation}
\Tr F_1
\le\Tr(A\rho_0^+A^\dagger).
\label{eq:proof-score-bound}
\end{equation}

Now suppose the range of $A$ lies in a projector $P_U$: $P_UA=A$. The positive operator on the left of Eq.~\eqref{eq:proof-operator-majorization} is then supported in $P_U$, hence
\begin{align}
\Tr F_1
&\le\Tr(A\rho_0^+A^\dagger)\\
&=\Tr[P_U(A\rho_0^+A^\dagger)]\\
&\le\frac4{R_{\rm lin}^2}\Tr(P_U\rho_0).
\end{align}
Thus
\begin{equation}
\frac{R_{\rm lin}^2}{4}\Tr F_1
\le\Tr(P_U\rho_0).
\label{eq:proof-upper-tail}
\end{equation}
The same argument applied to $A^\dagger$ yields a domain/lower-endpoint bound.

\paragraph{Finite-copy extension.}
For $N$ independently encoded copies,
\begin{equation}
A_N=\sum_{j=1}^N
\rho_0^{\otimes(j-1)}\otimes A\otimes
\rho_0^{\otimes(N-j)}.
\end{equation}
Because $\Tr A=0$, all cross-copy terms in
\begin{equation}
\Tr(A_N\rho_N^+A_N^\dagger)
\end{equation}
vanish, while every diagonal copy term equals the one-copy weighted norm. Therefore
\begin{equation}
\Tr(A_N\rho_N^+A_N^\dagger)
=N\Tr(A\rho_0^+A^\dagger).
\end{equation}
The weighted score inequality is valid for an arbitrary POVM on the full $N$-copy Hilbert space, so
\begin{equation}
\frac{R_{\rm lin}^2}{4}\frac{\Tr F_N}{N}
\le\Tr(P_U\rho_0).
\label{eq:proof-upper-tail-N}
\end{equation}
No product-measurement assumption is used.

\paragraph{Exact Bohr mode.}
If
\begin{equation}
[H,A_\nu]=\hbar\nu A_\nu,
\end{equation}
then the range and domain of $A_\nu$ lie in upper and lower endpoint sectors separated by $\hbar\nu$. For a stationary baseline, applying Eq.~\eqref{eq:proof-upper-tail-N} to $A_\nu$ and $A_\nu^\dagger$ gives
\begin{equation}
\frac{R_{\rm lin}^2}{4}\frac{\Tr F_N^{(\nu)}}{N}
\le\min\{U_\nu,D_\nu\}.
\end{equation}
For a semibounded Hamiltonian the upper endpoint lies inside the tail
\begin{equation}
T(\nu)=\Pr_{\rho_0}[H-E_*\ge\hbar\nu],
\end{equation}
so
\begin{equation}
\frac{R_{\rm lin}^2}{4}\frac{\Tr F_N^{(\nu)}}{N}
\le T(\nu).
\end{equation}
Markov's inequality then gives the mean-energy corollary.

\subsection{Autonomous dual survival}
\label{supp:dual-survival-full}

For the exact clock--signal exchange
\begin{equation}
[H_S,A_\nu]=+\hbar\nu A_\nu,
\qquad
[H_C,A_\nu]=-\hbar\nu A_\nu,
\end{equation}
the signal viewpoint sees $A_\nu$ as a positive $+\nu$ tangent, whereas the clock viewpoint sees $A_\nu^\dagger$ as a positive $+\nu$ tangent. The two Hermitian tangent quadratures span the same physical plane under $A_\nu\leftrightarrow A_\nu^\dagger$; only the sine coordinate reverses. Hence the affine radius and Fisher trace are unchanged.

Applying Eq.~\eqref{eq:proof-upper-tail-N} to the signal range projector and to the clock range projector of $A_\nu^\dagger$ yields
\begin{equation}
\frac{R_{\rm lin}^2}{4}\frac{\Tr F_N^{(\nu)}}{N}
\le T_S(\nu),
\end{equation}
and
\begin{equation}
\frac{R_{\rm lin}^2}{4}\frac{\Tr F_N^{(\nu)}}{N}
\le T_C(\nu).
\end{equation}
Therefore
\begin{equation}
\boxed{
\frac{R_{\rm lin}^2}{4}\frac{\Tr F_N^{(\nu)}}{N}
\le\min\{T_C(\nu),T_S(\nu)\}.
}
\end{equation}
The proof requires only global physicality of the baseline and exact local gap structure of the tangent; separate local stationarity is unnecessary because Eq.~\eqref{eq:proof-upper-tail-N} did not assume $[\rho_0,H_X]=0$.
